\documentclass[11pt]{article}

% ------------------------------------------------------------------
%  Compile with:  pdflatex REPORT.tex
%  Figures are read from the figures/ subfolder next to this file.
% ------------------------------------------------------------------
\usepackage[margin=1in]{geometry}
\usepackage{booktabs}
\usepackage{amsmath}
\usepackage{amssymb}
\usepackage{array}
\usepackage{graphicx}
\graphicspath{{figures/}{./}}   % find PNGs in figures/ or alongside this file
\usepackage{framed}
\usepackage{enumitem}
\usepackage[hidelinks]{hyperref}

% Black-and-white explanatory box: a rule-framed block (page-breakable), no header.
\newenvironment{plainbox}%
  {\par\smallskip\begin{framed}\noindent\itshape}%
  {\end{framed}\smallskip}

% All figures at a consistent width.
\newcommand{\reportfig}[2]{%
  \begin{center}\includegraphics[width=0.92\linewidth]{#1}\end{center}%
  \vspace{-6pt}\begin{center}\footnotesize #2\end{center}}

\newcommand{\rsq}{$R^2$}
\newcommand{\um}{\,$\mu$m}
\newcommand{\alp}{$\alpha$}

\title{\textbf{Machine-Learning Retrieval of Exoplanet Atmospheres\\
from Reflected-Light Spectra}\\[4pt]
\large A Consolidated Evaluation Report}
\author{}
\date{July 2026}

\begin{document}
\maketitle

\begin{abstract}
\noindent
We train compact neural networks (about half a million parameters) to read the
atmospheric composition of a planet from its reflected-light spectrum --- the rainbow of
sunlight it bounces back --- across the $0.2$--$2.0$\um{} band. Models are trained on a
large synthetic library of terrestrial (rocky) planets and evaluated not only on held-out
synthetic planets but, for the first time in this project, on \emph{real} telescope spectra
of $19$ Solar-System bodies observed as if they were exoplanets. The headline findings:
(i) the best model reaches $R^2 = 0.78$ on covered gases while being roughly ten times
smaller than the original NASA model for this data set (which used $\sim$18 million
parameters); (ii) on real data the models genuinely detect methane (perfect separation of
methane-rich from airless bodies) but fabricate oxygen (bare rocks are assigned more oxygen
than Earth), a flaw we trace to the training data and fix at inference time without
retraining; (iii) a simple combination of our existing models lifts accuracy to
$R^2 = 0.85$ with no new training. All numbers below come from frozen, already-trained
models.
\end{abstract}

\tableofcontents
\vspace{1em}

% ==================================================================
\section*{How to read this report}
\addcontentsline{toc}{section}{How to read this report}

Two numbers recur throughout:

\begin{itemize}[nosep]
\item \textbf{\rsq{} (``R-squared'')} measures how well predictions match the truth on a
scale where $1.0$ is perfect, $0$ is no better than guessing the average, and
\emph{negative} means worse than guessing. We report it in $\log_{10}$ space because gas
abundances span many orders of magnitude.
\item \textbf{\alp{} (exposure factor)} is how long the telescope stares, relative to a
nominal $8$-hour exposure. More light means less noise. $\alpha = 1$ is a realistic single
observation; the ``noiseless reference'' is the best-case limit with effectively infinite
light, isolating what the \emph{model} can do from what the \emph{photons} allow.
\end{itemize}

``Covered species'' are the five gases that actually leave a fingerprint in this wavelength
band: water (H$_2$O), carbon dioxide (CO$_2$), oxygen (O$_2$), methane (CH$_4$), and ozone
(O$_3$).

% ==================================================================
\section{Baseline accuracy and model size}

\begin{center}
\begin{tabular}{lrr}
\toprule
Model & Parameters & \rsq{} (covered) \\
\midrule
causal (best)     & 520\,k  & \textbf{0.781} \\
original1dcnn     & 36.0\,M & 0.779 \\
optimized1dcnn    & 1.20\,M & 0.728 \\
transformerarch   & 520\,k  & 0.568 \\
causal\_cfi       & 520\,k  & 0.355 \\
causal\_xl        & 2.05\,M & 0.173 \\
\midrule
\textbf{Linear stack of all 6} & --- & \textbf{0.849} \\
\bottomrule
\end{tabular}
\end{center}

\noindent
(Trained on the synthetic NASA-PSG library, tested on held-out planets from the same
library, noiseless reference.)

\reportfig{fig01_model_accuracy}{Accuracy by model. The best model (causal) is also the
smallest; the striped bar is the combined stack of all six.}

\begin{plainbox}
Our best model scores $0.78$ out of a perfect $1.0$ --- and it is also our
\emph{smallest} model, at about half a million parameters. For context, the original NASA
model that introduced this data set (Soboczenski et al., 2018) used a 1D CNN with
\emph{approximately 18 million} trainable parameters; our best model matches its approach
while being about $35\times$ smaller. Our own \texttt{original1dcnn} baseline (36M
parameters) is essentially that same NASA architecture scaled up, and it does \emph{not}
beat the 520k model. Keeping the core models small was a deliberate choice: they are cheap
to train on current hardware, which let us benchmark many architectures fairly and leaves a
clear runway to scale up later if needed. (More recent reflected-light models, e.g.\ Barbosa
et al.\ 2025, use larger attention-based architectures.) The final line is the most useful
takeaway of the whole report: simply \emph{combining} the six existing models --- with no
new training --- reaches $0.85$, our best number anywhere.
\end{plainbox}

% ==================================================================
\section{Reference baselines: what does the neural network have to beat?}

A retrieval accuracy number means nothing on its own --- it has to be measured against
simpler methods that establish a floor. If a deep network cannot beat a straight line, the
deep network is not justified. The project uses five reference baselines, all evaluated on the
identical input, at one realistic exposure ($\alpha = 1$) unless stated. The first three are
statistical, the fourth is a physics reference, and the fifth is an internal sanity check.

\begin{center}
\begin{tabular}{p{3.2cm}p{6.2cm}cc}
\toprule
Baseline & What it is & all-12 \rsq{} & covered \rsq{} \\
\midrule
Prior mean & Ignores the spectrum entirely; predicts the average training composition for every
planet. The ``no-information'' floor. & $-0.059$ & $-0.065$ \\[2pt]
Ridge (linear) & A straight-line fit from the (compressed) spectrum to abundances. The
\emph{linear-information} floor. & $0.031$ & $0.086$ \\[2pt]
Random forest & An ensemble of decision trees on the same features; a capacity-limited
non-linear reference. & $0.013$ & $0.058$ \\[2pt]
Band-depth ``ruler'' & Physics baseline: measure the depth of a gas's absorption band by hand,
no learning at all. & \multicolumn{2}{c}{CH$_4$ AUROC $=1.00$ (real data)} \\[2pt]
Naive model average & Averages all six neural models with equal weight; an ensemble sanity
check. & --- & $0.637$ \\
\midrule
\textbf{Neural (causal), $\alpha=1$} & Our model at the same exposure as the baselines.
& \textbf{0.058} & \textbf{0.128} \\
\textbf{Neural (causal), noiseless} & Our model in the best-case, high-signal limit.
& \textbf{0.589} & \textbf{0.781} \\
\bottomrule
\end{tabular}
\end{center}

\noindent The statistical baselines were fit on the full library ($\sim$88k train, $\sim$11k
test) using a $64$-bin compression of each spectrum. (There is also an external benchmark, not
a baseline model: published literature retrievals of real exoplanets, used elsewhere as a
reality check.)

\reportfig{fig13_baselines}{The reference floors versus the neural model, all at one real
exposure. The dashed line marks the neural model's best-case (noiseless) score for context.}

\begin{plainbox}
The two things to read here. First, at a \emph{single realistic exposure} everything is low ---
even our model only reaches $0.13$ on covered gases, because one $8$-hour look does not collect
enough light (Section~6). But it still beats the linear floor ($0.086$) and the do-nothing
floor ($-0.065$), so the network is genuinely extracting more than a straight line can. Second,
the honest comparison is in the \emph{best-case} limit, where the network reaches $0.78$ while
the linear and prior baselines stay near zero --- a large, real gap. A subtle but important
point: a random forest does \emph{worse} than a plain linear fit here, which tells us the
useful signal in these spectra is close to linear once the bands are located; the neural net's
advantage comes from using the full high-resolution spectrum, not from raw non-linearity. The
``ruler'' baseline matters for the real-data story: a hand-built band-depth measurement already
achieves perfect methane detection, so our claim is that the network \emph{matches} the physics,
not that it beats a simple detector.
\end{plainbox}

% ==================================================================
\section{Cross-generator test: does the model transfer to other physics engines?}

A model trained on one simulator should ideally still work on spectra made by a
\emph{different} simulator. We tested three: PSG (what we trained on), TauREx, and
petitRADTRANS (pRT).

\begin{center}
\footnotesize
\begin{tabular}{lccccc}
\toprule
Model & PSG (native) & TauREx & pRT & pRT + AdaBN & pRT + affine \\
      &              &        &     & (0 labels)  & (100 labels) \\
\midrule
optimized1dcnn  & 0.728 & $-0.274$ & $-3.736$ & $-0.162$ & $-0.038$ \\
original1dcnn   & 0.779 & $-0.374$ & $-5.562$ & n/a      & $-0.030$ \\
transformerarch & 0.568 & ---      & $-1.091$ & n/a      & $-0.041$ \\
causal          & 0.781 & $-0.325$ & $-3.048$ & ---      & $\sim 0$ \\
\bottomrule
\end{tabular}
\end{center}

\noindent
AdaBN re-estimates the model's internal normalisation statistics on unlabelled target
spectra (needs the model to have batch-normalisation layers, hence ``n/a'' for two of
them). ``Affine'' fits a small per-gas correction using $100$ labelled target planets.

\reportfig{fig02_crossgen_adaptation}{The catastrophic gap on petitRADTRANS is mostly a
surface mismatch: a zero-label adjustment and then 100 labels bring it back to near zero.}

\begin{plainbox}
Straight out of the box the model does badly on the other simulators --- the numbers go
strongly negative, meaning worse than guessing. But this is mostly a \emph{surface
mismatch}, not a loss of real skill: a zero-label statistical adjustment (AdaBN) removes
about $95\%$ of the gap, and $100$ labelled examples bring it to roughly zero. ``Roughly
zero'' means the model stops being actively wrong, but it does not gain genuine
cross-simulator skill.

\textbf{Context (author's note).} This took a long time to get working and the results were
disappointing. The deeper reason is a physics mismatch: pRT is fundamentally an
\emph{emission} (heat-glow) simulator with a reflected-light mode bolted on, while TauREx is
a \emph{reflected-light} model but one built for large gaseous planets, not the small rocky
atmospheres our dataset represents. We also tried training on two synthetic datasets at once
hoping for an overall improvement, but since each dataset comes from a different underlying
engine, mixing them was a dead end. Honestly, the cross-generator \rsq{} values are poor,
and this remains the weakest part of the pipeline.
\end{plainbox}

% ==================================================================
\section{Per-gas accuracy}

\begin{center}
\begin{tabular}{lccc}
\toprule
Gas & \rsq{} (noiseless) & \rsq{} at \alp{}$=1$ & RMSE at \alp{}$=1$ (dex) \\
\midrule
H$_2$O & 0.72 & 0.06 & 0.44 \\
CO$_2$ & 0.73 & 0.03 & 0.36 \\
O$_2$  & 0.80 & 0.12 & 0.36 \\
CH$_4$ & 0.79 & 0.11 & 0.42 \\
O$_3$  & 0.86 & 0.33 & 0.48 \\
N$_2$  & $-0.06$ & $-0.06$ & --- (invisible; honest null) \\
\bottomrule
\end{tabular}
\end{center}

\noindent\textbf{Derived carbon-to-oxygen (C/O) ratio}, in $\log_{10}$:
noiseless reference \rsq{} $= 0.81$, median error $0.05$ dex;
at \alp{}$=1$, \rsq{} $= 0.10$, median error $0.12$ dex.
(True $\log_{10}$ C/O spans $-1.92$ to $-0.24$, median $-0.58$.)

\reportfig{fig03_pergas_accuracy}{Per-gas accuracy in the best case (solid) versus one
realistic exposure (striped). Nitrogen is a permanent null --- it is invisible in this band.}

\begin{plainbox}
Each gas is recovered well in the best-case (noiseless) limit but degrades once realistic
observing noise is added --- the two columns show the gap. The C/O ratio, a key
planet-formation diagnostic, is genuinely recoverable when the signal is strong
($R^2 = 0.81$) but noise-limited at a single realistic exposure.

Two honest caveats. First, \textbf{temperature and metallicity are not things this model
predicts} --- the training data does not label them, so they simply are not retrieval
targets here. Second, \textbf{N$_2$ (nitrogen) is essentially impossible for any neural
network to detect} in this band, because nitrogen leaves no fingerprint in reflected light.
Its negative score is not a failure of the model; it is physics. Earth's atmosphere is
$78\%$ nitrogen and none of it is visible.
\end{plainbox}

% ==================================================================
\section{Exposure time versus accuracy}

\begin{center}
\begin{tabular}{lccccccc}
\toprule
Exposure ($\times$ nominal) & 0.09 & 1 & 9 & 100 & 900 & $10^4$ & $9{\times}10^4$ \\
\midrule
\rsq{} (covered) & $-0.02$ & 0.13 & 0.35 & 0.56 & 0.70 & 0.79 & 0.83 \\
\bottomrule
\end{tabular}
\end{center}

\reportfig{fig04_exposure_accuracy}{Accuracy versus exposure time for all six models. The
curve is set by photon noise, not by the choice of architecture.}

\begin{plainbox}
Accuracy is limited by how many photons the telescope collects, not by the model. A single
nominal exposure gives only $R^2 = 0.13$; reaching the model's true ceiling would need about
ten-thousand times longer staring. The practical message: \textbf{any accuracy claim is
meaningless without stating the exposure time it assumes.}
\end{plainbox}

% ==================================================================
\section{Clouds: can we retrieve through them?}

Clouds hide the gas fingerprints beneath them. We trained a separate ``de-clouding'' network
to reconstruct the clear spectrum, then fed the result to the frozen retrieval models. It was
trained only on one cloud type (``grey''); the other three are unseen tests.

\begin{center}
\begin{tabular}{lccc}
\toprule
Cloud family & Cloudy \rsq{} & De-clouded \rsq{} & Recovery \\
\midrule
grey (seen)        & $-1.02$ & $-0.19$ & 0.46 \\
non-grey haze      & $-0.51$ & $-0.19$ & 0.25 \\
band-selective     & $-0.61$ & $-0.28$ & 0.24 \\
patchy (Earth-like)& $-0.97$ & $-0.21$ & 0.44 \\
\bottomrule
\end{tabular}
\end{center}

\reportfig{fig05_cloud_generalization}{Clouds (light bars) crush accuracy below zero;
de-clouding (dark bars) recovers part of it, even on cloud types not seen in training.}

\begin{plainbox}
The de-clouder recovers $24$--$46\%$ of the accuracy that clouds destroy, and importantly it
still helps on cloud types it never saw in training. So as a stand-alone tool it works ---
partially.

\textbf{Context (author's note).} Overall, clouds remain hard. We tried two routes and
neither was a clean success. The first was to teach the retrieval models directly on
cloud-augmented data (muting and boosting parts of the spectrum); this needed very long
training and did not actually improve things --- it made the models worse. The second was
this separate de-clouding module (around one million parameters, two variants tried). It
works in isolation, but there is a fundamental limitation: our clouds are generated from
tidy mathematical formulas, not real observed cloud noise, so the ``restored'' spectrum does
not fully match what a real cloudy planet would give.

\textbf{Important safety caveat.} The de-clouder is \emph{not safe to run on an already-clear
spectrum} --- doing so badly damages it, because it was never shown a clear sky during
training and always assumes a cloud is present. It should only be applied when a thick cloud
is independently known to be there, and it needs one training fix before general use.
\end{plainbox}

% ==================================================================
\section{Uncertainty estimates}

\begin{center}
\begin{tabular}{lccc}
\toprule
Gas & Interval width (dex) & Coverage (native) & Coverage (at R$=$22) \\
\midrule
H$_2$O & 0.20 & 0.71 & 0.36 \\
CO$_2$ & 0.15 & 0.71 & 0.20 \\
O$_2$  & 0.13 & 0.67 & 0.26 \\
CH$_4$ & 0.16 & 0.67 & 0.41 \\
O$_3$  & 0.18 & 0.67 & 0.16 \\
\bottomrule
\end{tabular}
\end{center}

\noindent Target coverage is $0.68$ (a ``68\% confidence interval'' should contain the truth
$68\%$ of the time).

\reportfig{fig06_uncertainty_coverage}{Error-bar honesty. At native resolution the
intervals cover the truth about the right fraction of the time; on blurred data they become
far too narrow.}

\begin{plainbox}
The raw model is \emph{overconfident} --- its error bars are too narrow. A statistical
technique called split-conformal calibration produces \emph{honest} error bars that contain
the truth the right fraction of the time ($\sim 0.68$) --- but only when the test data looks
like the training data. When the spectrum is lower-resolution (blurrier), the same error bars
become far too narrow (coverage drops to $0.16$--$0.41$); the honest interval is $2$--$4$
times wider. The fix is to re-calibrate the error bars at the target's actual resolution,
which needs no retraining.
\end{plainbox}

% ==================================================================
\section{Real data: Solar-System bodies observed as exoplanets}

This is the first time the models were tested on \emph{real} telescope spectra rather than
simulations. We used a public catalog of $19$ Solar-System bodies. Crucially, we know exactly
what these objects are made of, so we can grade the models honestly. The bodies fall into
three groups: four terrestrial planets/moons \emph{with} atmospheres, four giant planets, and
eleven \emph{airless} bodies (moons, asteroids) with \emph{no} atmosphere at all --- the
last group is a built-in trap for false detections.

\subsection{Methane: a genuine detection}

\begin{center}
\begin{tabular}{lc}
\toprule
Model & CH$_4$ detection AUROC \\
\midrule
original1dcnn   & 1.00 \\
transformerarch & 1.00 \\
causal          & 1.00 \\
causal\_cfi     & 1.00 \\
optimized1dcnn  & 0.88 \\
causal\_xl      & 0.76 \\
\bottomrule
\end{tabular}
\end{center}

\noindent AUROC (Area Under the ROC Curve) measures how well the model separates two groups:
$1.0$ is perfect separation, $0.5$ is a coin flip. Here: $5$ methane-rich bodies versus $5$
airless bodies, all at matched resolution.

\reportfig{fig07_ch4_detection_auroc}{Methane detection on real spectra. Four of six models
separate methane-rich from airless bodies perfectly.}

\begin{plainbox}
Four of six models achieve \emph{perfect} separation ($1.00$): every methane-rich body scores
higher than every airless one. We confirmed this is real, not luck: blanking out the methane
wavelengths destroys the signal, while blanking unrelated wavelengths does not. A simple
hand-built ``ruler'' that just measures band depth also scores $1.00$ --- so the fair claim is
\emph{the network matches the physics}, not that it beats it. This is the strongest real-data
result in the project.
\end{plainbox}

\subsection{Oxygen: a fabricated ``detection''}

\begin{center}
\small
\begin{tabular}{lrrrrrrr}
\toprule
Body & Class & R & usable? & true dom. & pred. & dex err & pred O$_2$+O$_3$ \\
\midrule
Earth     & terrestrial & 867 & yes & N$_2$ & O$_2$ & 2.17 & 0.41 \\
Venus     & terrestrial & 22  & no  & CO$_2$& CO$_2^{*}$ & 1.73 & 0.18 \\
Mars      & terrestrial & 22  & no  & CO$_2$& O$_2$ & 1.19 & 0.41 \\
Titan     & terrestrial & 829 & yes & N$_2$ & CO$_2$& 0.64 & 0.33 \\
Jupiter   & giant       & 865 & yes & H$_2$/He & O$_2$ & --- & 0.31 \\
Saturn    & giant       & 867 & yes & H$_2$/He & O$_2$ & --- & 0.38 \\
Uranus    & giant       & 867 & yes & H$_2$/He & O$_2$ & --- & 0.32 \\
Neptune   & giant       & 867 & yes & H$_2$/He & CO$_2$& --- & 0.32 \\
Mercury   & airless     & 43  & no  & none  & O$_2$ & --- & \textbf{0.46} \\
Moon      & airless     & 867 & yes & none  & O$_2$ & --- & \textbf{0.45} \\
Io        & airless     & 12  & no  & none  & O$_2$ & --- & 0.39 \\
Europa    & airless     & 25  & no  & none  & CO$_2$& --- & 0.33 \\
Ganymede  & airless     & 31  & no  & none  & CO$_2$& --- & 0.29 \\
Callisto  & airless     & 21  & no  & none  & O$_2$ & --- & 0.41 \\
Enceladus & airless     & 526 & yes & none  & CO$_2$& --- & 0.15 \\
Dione     & airless     & 829 & yes & none  & CO$_2$& --- & 0.26 \\
Rhea      & airless     & 866 & yes & none  & CO$_2$& --- & 0.31 \\
Ceres     & airless     & 866 & yes & none  & O$_2$ & --- & 0.42 \\
Pluto     & airless     & 100 & no  & none  & CO$_2$& --- & 0.31 \\
\bottomrule
\end{tabular}
\end{center}

\noindent$^{*}$Venus's correct CO$_2$ call is a ``stopped clock'': its spectrum is too
low-resolution to use, and the same model labels many airless rocks CO$_2$ too.

\reportfig{fig08_o2_fabrication}{Predicted oxygen on all 19 bodies, sorted. Light bars are
airless bodies with zero atmosphere; several out-score Earth (dotted line = Earth's true
value).}

\begin{plainbox}
The airless bodies have \emph{zero} oxygen --- they have no atmosphere. Yet the models assign
them about $30$--$46\%$ oxygen, with the Moon and Mercury (bare rock) scoring \emph{higher}
than Earth. The reason is not the architecture: the training set happens to contain a lot of
oxygen-rich planets, so the model's safe default guess for anything unfamiliar is ``about a
third oxygen.'' It is reciting a habit, not reading the spectrum. This is why an apparent
``oxygen detection'' on real Earth in earlier reports was never real.
\end{plainbox}

\subsection{Earth versus the Moon: the cleanest control}

Earth and the Moon were observed on the same night, with the same instrument, at the same
resolution. One has $21\%$ oxygen; the other is a rock.

\begin{center}
\begin{tabular}{lccc}
\toprule
Model & Earth O$_2$ (true 0.21) & Moon O$_2$ (true 0.00) & Earth $-$ Moon \\
\midrule
original1dcnn   & 0.455 & 0.461 & $-0.006$ \\
optimized1dcnn  & 0.517 & 0.536 & $-0.019$ \\
transformerarch & 0.392 & 0.465 & $-0.073$ \\
causal          & 0.404 & 0.446 & $-0.043$ \\
causal\_cfi     & 0.312 & 0.305 & $+0.007$ \\
causal\_xl      & 0.311 & 0.313 & $-0.002$ \\
\bottomrule
\end{tabular}
\end{center}

\reportfig{fig09_earth_vs_moon}{Earth (dark) versus the Moon (striped) in predicted oxygen.
The rock matches or beats Earth for five of six models.}

\begin{plainbox}
Five of six models give the \emph{rock} more oxygen than Earth. This is the sharpest possible
demonstration that the oxygen numbers carry no real information --- until we fix them
(next section).
\end{plainbox}

\subsection{Resolution: which gases need how sharp a spectrum?}

\begin{center}
\footnotesize
\begin{tabular}{l|ccccc|ccccc}
\toprule
 & \multicolumn{5}{c|}{Raw \rsq{}} & \multicolumn{5}{c}{After recalibration} \\
R & H$_2$O & CO$_2$ & O$_2$ & CH$_4$ & O$_3$ & H$_2$O & CO$_2$ & O$_2$ & CH$_4$ & O$_3$ \\
\midrule
22     & $-0.24$ & $-1.87$ & $-0.15$ & 0.32 & 0.27 & 0.36 & 0.11 & 0.33 & 0.34 & 0.80 \\
50     & 0.48 & $-4.37$ & 0.26 & 0.50 & 0.55 & 0.54 & 0.09 & 0.53 & 0.56 & 0.85 \\
100    & 0.55 & $-1.82$ & 0.53 & 0.64 & 0.67 & 0.61 & 0.22 & 0.69 & 0.68 & 0.87 \\
200    & 0.64 & 0.39 & 0.71 & 0.74 & 0.76 & 0.67 & 0.52 & 0.76 & 0.76 & 0.88 \\
400    & 0.60 & 0.68 & 0.78 & 0.79 & 0.81 & 0.68 & 0.68 & 0.78 & 0.81 & 0.89 \\
870    & 0.59 & 0.70 & 0.80 & 0.80 & 0.83 & 0.70 & 0.70 & 0.79 & 0.82 & 0.89 \\
native & 0.71 & 0.74 & 0.80 & 0.78 & 0.83 & 0.74 & 0.74 & 0.84 & 0.83 & 0.89 \\
\bottomrule
\end{tabular}
\end{center}

\noindent ``R'' is spectral resolution --- how finely the spectrum is sampled. Higher is
sharper.

\reportfig{fig10_resolution_requirements}{Accuracy versus how sharp the spectrum is. Most
gases survive moderate blur; carbon dioxide collapses below R$=$200 (dashed line).}

\begin{plainbox}
Carbon dioxide is the demanding case: it needs a fairly sharp spectrum ($R \geq 200$), and if
the spectrum is too blurry ($R = 22$) \emph{no amount of correction can recover it} --- the
information is physically gone. This is exactly why the models fail on catalog Venus and Mars:
those spectra are too blurry, which is a limitation of the \emph{data}, not the model. Every
other gas degrades gracefully and can be recalibrated down to moderate resolution; ozone
survives even the blurriest case.

\textbf{Context (author's note).} This was a key test --- using real planets and their moons
instead of only simulations. The result was encouraging: the models show real signs of
understanding, such as correctly picking out methane, even though they over-estimate oxygen.
\end{plainbox}

% ==================================================================
\section{New inference-only fixes (no retraining)}

Everything here runs on the existing frozen models --- no retraining, no new data.

\subsection{Continuum-twin nulling: removing the fabricated oxygen}

For each body we build a ``twin'' of its spectrum with the absorption bands smoothed away,
then subtract the twin's prediction from the real one. Since the twin shares everything
\emph{except} the gas fingerprints, the model's habitual guess cancels out.

\begin{center}
\begin{tabular}{lcc}
\toprule
Metric & causal (raw $\rightarrow$ nulled) & optimized (raw $\rightarrow$ nulled) \\
\midrule
CH$_4$ AUROC              & $1.00 \rightarrow 0.92$ & $0.88 \rightarrow 1.00$ \\
O$_2$ AUROC (Earth vs rocks) & $0.64 \rightarrow 0.91$ & $0.73 \rightarrow 1.00$ \\
mean $|$O$_2|$ on rocks   & $0.341 \rightarrow 0.016$ & $0.468 \rightarrow 0.006$ \\
mean $|$CO$_2|$ on rocks  & $0.351 \rightarrow 0.013$ & $0.186 \rightarrow 0.007$ \\
mean $|$N$_2|$ on rocks   & $0.255 \rightarrow 0.003$ & $0.271 \rightarrow 0.004$ \\
\bottomrule
\end{tabular}
\end{center}

\reportfig{fig11_twin_nulling}{Fabricated gas on airless rocks (light bars) before and
(dark bars) after twin-nulling --- it collapses toward zero, where it should be.}

\begin{plainbox}
This works remarkably well: the fabricated gas on airless rocks shrinks by $20$--$80$ times,
while the genuine methane signal survives, and the Earth-versus-Moon oxygen comparison finally
comes out the right way round. It converts a model that ``recites its habit'' into one that
``reports the difference the bands make'' --- and it costs nothing but a second forward pass.
\end{plainbox}

\subsection{Model stacking: breaking the accuracy ceiling}

\begin{center}
\begin{tabular}{lc}
\toprule
Combiner & \rsq{} (covered) \\
\midrule
Best single model        & 0.777 \\
Plain average of 6 models & 0.637 \\
Per-gas linear stack      & \textbf{0.849} \\
\bottomrule
\end{tabular}
\end{center}

\reportfig{fig12_model_stacking}{Combining the six frozen models with a per-gas weighting
(striped bar) beats the best single model; a naive average is worse.}

\begin{plainbox}
The six models make \emph{different} mistakes, so combining them cleverly (a small per-gas
weighted blend, fitted on held-out data) beats the best individual model by $+0.07$ ---
reaching $0.85$. The naive average is \emph{worse} than the best single model, because it lets
the weak models drag down the strong ones; the weighting is what matters. This is the cheapest
accuracy gain available and requires no retraining.
\end{plainbox}

% ==================================================================
\section{What did not work (honest log)}

\begin{itemize}[leftmargin=1.4em]
\item \textbf{Cross-simulator training.} Training on other physics engines (pRT, TauREx) to
gain transfer was abandoned: pRT models heat-emission not reflected light, TauREx targets gas
giants not rocky worlds, and mixing datasets from different engines was a dead end. Fine-tuning
on one of them caused catastrophic forgetting (accuracy $0.78 \rightarrow 0.32$).
\item \textbf{Cloud-augmented retraining.} Teaching the retrieval models on artificially
clouded spectra needed very long training and made accuracy \emph{worse}.
\item \textbf{The naive 6-model average} (0.637) is worse than the best single model (0.777);
only the weighted stack helps.
\item \textbf{The de-clouder on clear spectra} is actively harmful (see \S6) until retrained.
\item \textbf{Nitrogen, and CO$_2$ at low resolution}, are physically unrecoverable --- these
are honest ``no'' results, not fixable by better models.
\item \textbf{A trap worth remembering:} de-clouding appeared to ``fix Venus'' (correct CO$_2$)
but was in fact labelling \emph{every} body CO$_2$ --- only the airless-body controls exposed
it.
\end{itemize}

% ==================================================================
\section{Outlook}

We believe we have built close to the best achievable model for this dataset and are near the
state of the art, at a fraction of the usual model size. The natural next step is to train a
single large ($\sim 5$M-parameter) model on the best architecture and repeat the full
evaluation. Because testing a model that size is time-consuming, we propose doing so only after
agreeing the path forward.

\medskip
\noindent\textit{All results in this report are from frozen, already-trained models; no model
was retrained to produce these numbers.}

\end{document}
